\section{Validation scheme, tuning grids, random seeds, and computing environment}
\label{sec:validation}

\begin{itemize}
  \item Global seed: 42 (splits, CV shuffling, models, resampling, corruption, bootstrap).
  \item Hyper-parameter selection: \texttt{StratifiedKFold(5, shuffle=True, random\_state=42)} on the training
        partition only (regression: \texttt{KFold(5)}). Selection metric: Macro-F1 (classification), RMSE (regression).
        The selected configuration is refit on the full training partition and evaluated on validation.
  \item The test partitions are read only by \texttt{src/q5\_final\_test.py}; the loaders raise an error when
        called from anywhere else, and a guard file prevents a second evaluation.
  \item Classification data were \emph{not} subsampled: all \NTotal{} rows are used by all models.
\end{itemize}

\paragraph{Search spaces.} (from \texttt{src/config.py})
\begin{itemize}
  \item LogReg: $C \in \{0.001, 0.01, 0.1, 1, 10, 100\}$ $\times$ \texttt{class\_weight} $\in \{\text{None}, \text{balanced}\}$.
  \item Naive Bayes: GaussianNB \texttt{var\_smoothing} $\in \{10^{-9}, 10^{-6}, 10^{-3}, 10^{-1}\}$;
        BernoulliNB $\alpha \in \{0.01, 0.1, 1, 10\}$ $\times$ number of quantile bins $\in \{3, 5, 10\}$.
  \item $k$-NN: $k \in \{1, 3, 5, 9, 15, 31, 51\}$ $\times$ weights $\in$ \{uniform, distance\} $\times$ $p \in \{1, 2\}$.
  \item Decision Tree: criterion $\in$ \{gini, entropy\} $\times$ max\_depth $\in \{3, 5, 8, 12, \text{None}\}$ $\times$
        min\_samples\_leaf $\in \{1, 5, 20, 50\}$ $\times$ ccp\_alpha $\in \{0, 10^{-4}, 10^{-3}\}$ $\times$
        class\_weight $\in$ \{None, balanced\} (log\_loss checked separately).
  \item Random Forest (100 trees in the grid; 60 trees for the depth/max\_features curves): max\_depth $\in \{12, \text{None}\}$ $\times$ max\_features $\in$
        \{sqrt, 0.5\} $\times$ min\_samples\_leaf $\in \{1, 5\}$ $\times$ class\_weight $\in$ \{None,
        balanced\_subsample\}; curves over n\_estimators $\in \{10, \dots, 400\}$, max\_depth and max\_features.
  \item AdaBoost: learning rate $\in \{0.05, 0.1, 0.5, 1\}$; GradientBoosting: learning rate $\in \{0.05, 0.1, 0.2\}$;
        both $\times$ balanced sample weights $\in$ \{off, on\} $\times$ n\_estimators $\in \{10, 25, 50, 100, 150, 200, 300, 400\}$.
  \item Regression: Ridge $\alpha \in$ logspace$(-4, 3, 36)$; Lasso $\alpha \in$ logspace$(-5, 0, 31)$; tree max\_depth
        $\in \{4, 6, 8, 10, 12, 15, \text{None}\}$ $\times$ min\_samples\_leaf $\in \{1, 5, 20\}$; RF (200 trees)
        max\_depth $\in \{10, 20, \text{None}\}$ $\times$ max\_features $\in \{0.33, 1.0\}$ $\times$ min\_samples\_leaf $\in \{1, 5\}$.
\end{itemize}
\inputtable{q2_best_params}
\inputtable{q4_best_params}
\inputtable{q3_best_params}

For classification, Macro-F1 is selected as the primary optimization and selection metric because accuracy is severely misleading under an 11.7\% positive class prevalence: a majority-class Dummy classifier achieves 0.883 accuracy while identifying zero subscribers (recall $= 0.000$). Selecting by accuracy would reward models that ignore the minority class. Macro-F1 averages the $F_1$ scores of the \texttt{no} and \texttt{yes} classes with equal weight, immediately exposing the Dummy baseline's failure (0.469) and requiring balanced performance on both classes. Although ROC-AUC and PR-AUC are valuable threshold-independent ranking diagnostics and are reported alongside Macro-F1, deploying a telemarketing campaign ultimately requires a concrete binary call/no-call decision, making a thresholded F1 metric more directly aligned with operational decisions. For regression on California Housing, RMSE is chosen over MAE because squaring residuals penalizes large valuation errors disproportionately; in real-estate pricing, a handful of \$200,000--\$300,000 mispredictions carry far greater financial risk than many small errors, and RMSE remains interpretable in the target's \$100,000 units while matching standard squared-error objectives.

\paragraph{Computing environment.} Generated by \texttt{src/environment.py}:
\IfFileExists{tables/environment.tex}{\begin{verbatim}
OS: Windows-10-10.0.26300-SP0
CPU: Intel64 Family 6 Model 165 Stepping 2, GenuineIntel | physical cores: 4 | logical cores: 8
RAM: 15.9 GiB
Python: 3.11.2
numpy: 2.1.3
scipy: 1.14.1
pandas: 2.2.3
scikit-learn: 1.5.2
imbalanced-learn: 0.12.4
matplotlib: 3.9.2
seaborn: 0.13.2
ucimlrepo: 0.0.7
joblib: 1.6.0
SEED: 42 | CV folds: 5 | n_jobs: -1
\end{verbatim}
}{}
